\documentclass[11pt,a4paper]{article}
\usepackage{layout}
\begin{document}

% Page 1: restore the approved 2026 chapter opener.
\begin{canvas}
\pdfbookmark[0]{第一章 从开放协作到公共价值}{chapter1}
\node[anchor=north west,inner sep=0pt,outer sep=0pt] at (current page.north west){\includegraphics[width=210mm,height=297mm]{.build/assets/chapter1-cover.pdf}};
\Box{27}{227}{161}{\sffamily\bfseries\fontsize{11.5}{16}\selectfont\color{teal}AI 重塑开源，开源塑造 AI}
\end{canvas}

% Page 2: AI -> open collaboration, with distinct external windows.
\begin{canvas}
\Header{1.1 · 双向重塑}{blue}\Footer{02}
\Sec{blue}{1.1　人工智能时代的开源协作：双向重塑}{section11}
\Title{AI 改变创造与协作的节奏}
\Subhead{59}{blue}{从活动增长，进入生产方式与协作责任的变化}
\Body{22}{72}{166}{AI 生成与智能体（可分解任务、调用工具的 AI 系统）进入开发过程，也扩展了模型与数据供给。GitHub（代码协作平台）和 Hugging Face（HF，模型、数据及应用平台）的官方报告呈现出不同侧面的增长。关键问题随之出现：谁创造，谁审阅，谁维护？}
\Label{22}{97}{166}{开发活动：GitHub 的三个年度观察}
\Box{22}{106}{52}{\sffamily\bfseries\fontsize{10}{15}\selectfont\color{blue}代码提交}
\Box{22}{114}{52}{\sffamily\fontsize{7.5}{10}\selectfont\color{muted}单位：十亿次}
\Rule{22}{160}{69}{line}
\Rect{25}{151.2536443148688}{10}{8.746355685131196}{blue!22}
\Box{22}{144.2536443148688}{16}{\sffamily\bfseries\fontsize{8.7}{12}\selectfont\color{navy}\centering 10.0}
\Box{22}{162}{16}{\sffamily\fontsize{7.8}{10}\selectfont\color{muted}\centering 2024}
\Rect{40}{148.71720116618076}{10}{11.282798833819243}{blue!47}
\Box{37}{141.71720116618076}{16}{\sffamily\bfseries\fontsize{8.7}{12}\selectfont\color{navy}\centering 12.9}
\Box{37}{162}{16}{\sffamily\fontsize{7.8}{10}\selectfont\color{muted}\centering 2025}
\Rect{55}{130.0}{10}{30.000000000000004}{blue}
\Box{52}{123.0}{16}{\sffamily\bfseries\fontsize{8.7}{12}\selectfont\color{navy}\centering 34.3}
\Box{52}{162}{16}{\sffamily\fontsize{7.8}{10}\selectfont\color{muted}\centering 2026}
\Box{79}{106}{52}{\sffamily\bfseries\fontsize{10}{15}\selectfont\color{blue}自动化作业}
\Box{79}{114}{52}{\sffamily\fontsize{7.5}{10}\selectfont\color{muted}单位：十亿次}
\Rule{79}{160}{126}{line}
\Rect{82}{154.7136563876652}{10}{5.286343612334802}{blue!22}
\Box{79}{147.7136563876652}{16}{\sffamily\bfseries\fontsize{8.7}{12}\selectfont\color{navy}\centering 8.0}
\Box{79}{162}{16}{\sffamily\fontsize{7.8}{10}\selectfont\color{muted}\centering 2024}
\Rect{97}{150.7488986784141}{10}{9.251101321585903}{blue!47}
\Box{94}{143.7488986784141}{16}{\sffamily\bfseries\fontsize{8.7}{12}\selectfont\color{navy}\centering 14.0}
\Box{94}{162}{16}{\sffamily\fontsize{7.8}{10}\selectfont\color{muted}\centering 2025}
\Rect{112}{130.0}{10}{30.0}{blue}
\Box{109}{123.0}{16}{\sffamily\bfseries\fontsize{8.7}{12}\selectfont\color{navy}\centering 45.4}
\Box{109}{162}{16}{\sffamily\fontsize{7.8}{10}\selectfont\color{muted}\centering 2026}
\Box{136}{106}{52}{\sffamily\bfseries\fontsize{10}{15}\selectfont\color{blue}合入修改提案}
\Box{136}{114}{52}{\sffamily\fontsize{7.5}{10}\selectfont\color{muted}单位：十亿次}
\Rule{136}{160}{183}{line}
\Rect{139}{150.72307692307692}{10}{9.276923076923078}{blue!22}
\Box{136}{143.72307692307692}{16}{\sffamily\bfseries\fontsize{8.7}{12}\selectfont\color{navy}\centering 0.402}
\Box{136}{162}{16}{\sffamily\fontsize{7.8}{10}\selectfont\color{muted}\centering 2024}
\Rect{154}{147.58461538461538}{10}{12.415384615384616}{blue!47}
\Box{151}{140.58461538461538}{16}{\sffamily\bfseries\fontsize{8.7}{12}\selectfont\color{navy}\centering 0.538}
\Box{151}{162}{16}{\sffamily\fontsize{7.8}{10}\selectfont\color{muted}\centering 2025}
\Rect{169}{130.0}{10}{30.0}{blue}
\Box{166}{123.0}{16}{\sffamily\bfseries\fontsize{8.7}{12}\selectfont\color{navy}\centering 1.30}
\Box{166}{162}{16}{\sffamily\fontsize{7.8}{10}\selectfont\color{muted}\centering 2026}
\Box{22}{183}{166}{\sffamily\bfseries\fontsize{10.5}{15}\selectfont\color{teal}HF 开放资产：模型、数据与应用空间同步扩张}
\Box{22}{195}{52}{\sffamily\bfseries\fontsize{9.5}{14}\selectfont\color{navy}公共模型仓库}
\Box{54}{195}{19}{\sffamily\bfseries\fontsize{9.5}{14}\selectfont\color{teal}\raggedleft +22\%}
\Rect{22}{208}{29.554054054054053}{4}{teal!25}
\Box{53.55405405405405}{206.5}{14}{\sffamily\fontsize{8.5}{11}\selectfont\color{navy}243}
\Rect{22}{218}{36.0}{4}{teal}
\Box{60.0}{216.5}{14}{\sffamily\fontsize{8.5}{11}\selectfont\color{navy}296}
\Box{79}{195}{52}{\sffamily\bfseries\fontsize{9.5}{14}\selectfont\color{navy}数据集}
\Box{111}{195}{19}{\sffamily\bfseries\fontsize{9.5}{14}\selectfont\color{teal}\raggedleft +41\%}
\Rect{79}{208}{25.596}{4}{teal!25}
\Box{106.596}{206.5}{14}{\sffamily\fontsize{8.5}{11}\selectfont\color{navy}71.1}
\Rect{79}{218}{36.0}{4}{teal}
\Box{117.0}{216.5}{14}{\sffamily\fontsize{8.5}{11}\selectfont\color{navy}100}
\Box{136}{195}{52}{\sffamily\bfseries\fontsize{9.5}{14}\selectfont\color{navy}应用空间}
\Box{168}{195}{19}{\sffamily\bfseries\fontsize{9.5}{14}\selectfont\color{teal}\raggedleft +44\%}
\Rect{136}{208}{25.0}{4}{teal!25}
\Box{163.0}{206.5}{14}{\sffamily\fontsize{8.5}{11}\selectfont\color{navy}100}
\Rect{136}{218}{36.0}{4}{teal}
\Box{174.0}{216.5}{14}{\sffamily\fontsize{8.5}{11}\selectfont\color{navy}144}
\Box{22}{228}{166}{\sffamily\fontsize{7.8}{11}\selectfont\color{muted}浅色为来源所列期初，深色为期末；单位：万个。增幅按展示数值计算。}

\Body{22}{239}{166}{PR（合并请求）是提交修改、请他人审阅并合入的提案。GitHub 观察到，2026 年 9 月逾 40\% 的 PR 借助智能体发起。自动化作业执行测试等任务；应用空间托管 AI 演示与应用。}
\Note{22}{260}{166}{图 1-1　来源：\href{\GH}{GitHub Octoverse 2026}、\href{\HF}{HF Summer 2026}。GitHub 各年为截至 9 月 30 日的 12 个月，含公开与私有仓库；HF 为夏季报告所列公共资源存量，具体快照日未明示。各面板独立标度；两平台不合并计算。}
\end{canvas}

% Page 3: complementary evidence of reuse, usage and cross-border continuation.
\begin{canvas}
\Header{1.1 · 双向重塑}{blue}\Footer{03}
\Kicker{blue}{1.1 / OPEN COLLABORATION SHAPES AI}
\Title{开源改变 AI 的扩散与演进}
\Subhead{59}{blue}{中国模型与全球社区：开放供给如何形成全球接续}
\Body{22}{72}{166}{开放供给的影响还包括全球社区的再加工与使用。衍生仓库记录基于已有模型的继续开发；推理指用模型生成结果，词元是模型处理文本的基本单元。Qwen 与 DeepSeek 展现了不同的生态位置。}
\shade[left color=blue!10,right color=blue!2,rounded corners=2mm,draw=blue!17,line width=.45pt] \Pos{22}{99} rectangle \Pos{102}{163};
\path[fill=blue!9] \Pos{50.8}{163} .. controls \Pos{76.4}{162} and \Pos{82.8}{158} .. \Pos{101}{159} -- \Pos{101}{162} -- \Pos{50.8}{163} -- cycle;
\Box{26}{103}{72}{\sffamily\bfseries\fontsize{10.3}{15}\selectfont\color{blue}Qwen｜广泛的社区派生}
\shade[left color=teal!10,right color=teal!2,rounded corners=2mm,draw=teal!17,line width=.45pt] \Pos{108}{99} rectangle \Pos{188}{163};
\path[fill=teal!9] \Pos{136.8}{163} .. controls \Pos{162.4}{162} and \Pos{168.8}{158} .. \Pos{187}{159} -- \Pos{187}{162} -- \Pos{136.8}{163} -- cycle;
\Box{112}{103}{72}{\sffamily\bfseries\fontsize{10.3}{15}\selectfont\color{teal}DeepSeek｜大型模型与推理使用}
\Box{26}{115}{72}{\sffamily\bfseries\fontsize{24}{28}\selectfont\color{blue}151,448}
\Box{26}{129}{72}{\sffamily\fontsize{9}{13}\selectfont\color{navy}HF 报告记录的社区衍生仓库}
\Box{26}{141}{72}{\sffamily\fontsize{8.4}{12}\selectfont\color{muted}GGUF 是便于本地运行的模型文件格式。}
\Box{26}{149}{72}{\sffamily\fontsize{8.4}{12}\selectfont\color{muted}28,531 个转换中，官方发布 54 个；约 99.8\% 来自其他账号。}
\Box{112}{116}{34}{\sffamily\bfseries\fontsize{23}{28}\selectfont\color{teal}47\%}
\Box{151}{116}{33}{\sffamily\bfseries\fontsize{23}{28}\selectfont\color{teal}31.1\%}
\Box{112}{129}{34}{\sffamily\fontsize{8.5}{12}\selectfont\color{navy}所追踪模型下载份额}
\Box{151}{129}{33}{\sffamily\fontsize{8.5}{12}\selectfont\color{navy}所选样本词元份额}
\Box{112}{141}{72}{\sffamily\fontsize{8.4}{12}\selectfont\color{muted}总参数 2500 亿及以上模型；}
\Box{112}{149}{72}{\sffamily\fontsize{8.4}{12}\selectfont\color{muted}OpenRouter 模型调用平台的前十开放模型样本。}
\Box{22}{169}{166}{\sffamily\bfseries\fontsize{10.5}{15}\selectfont\color{navy}从模型输出，到开放数据，再到后续模型}
\draw[teal!50,line width=1pt,-{Stealth[length=1.4mm]}] \Pos{53}{196} .. controls \Pos{61}{193} and \Pos{64}{199} .. \Pos{72}{196};
\draw[gold!50,line width=1pt,-{Stealth[length=1.4mm]}] \Pos{137}{196} .. controls \Pos{145}{193} and \Pos{149}{199} .. \Pos{157}{196};
\fill[blue!5] \Pos{41}{196} circle[radius=11.2mm];
\shade[top color=white,bottom color=blue!28,shading angle=30,draw=blue!35,line width=.45pt] \Pos{41}{196} circle[radius=10mm];
\draw[white,line width=.7pt] \Pos{35.2}{190.2} arc[start angle=135,end angle=65,radius=8.2mm];
\Box{35.8}{190.8}{10.4}{\sffamily\fontsize{9.5}{14}\selectfont\color{ink}\includegraphics[width=10.4mm]{.build/assets/D04.pdf}}
\Box{22}{209}{39}{\sffamily\bfseries\fontsize{9.8}{13}\selectfont\color{blue}\centering DeepSeek-R1}
\shade[left color=teal!10,right color=teal!2,rounded corners=2mm,draw=teal!17,line width=.45pt] \Pos{75}{181} rectangle \Pos{137}{217};
\path[fill=teal!9] \Pos{97.32}{217} .. controls \Pos{117.16}{216} and \Pos{122.12}{212} .. \Pos{136}{213} -- \Pos{136}{216} -- \Pos{97.32}{217} -- cycle;
\Box{79}{184}{54}{\sffamily\bfseries\fontsize{9.1}{13}\selectfont\color{teal}NVIDIA 开放后训练数据}
\Box{79}{194}{54}{\sffamily\bfseries\fontsize{19}{23}\selectfont\color{teal}121.3 万条}
\Box{79}{207}{54}{\sffamily\fontsize{7.7}{11}\selectfont\color{muted}明确列为 R1 生成的样本}
\fill[gold!5] \Pos{169}{196} circle[radius=11.2mm];
\shade[top color=white,bottom color=gold!28,shading angle=30,draw=gold!35,line width=.45pt] \Pos{169}{196} circle[radius=10mm];
\draw[white,line width=.7pt] \Pos{163.2}{190.2} arc[start angle=135,end angle=65,radius=8.2mm];
\Box{163.8}{190.8}{10.4}{\sffamily\fontsize{9.5}{14}\selectfont\color{ink}\includegraphics[width=10.4mm]{.build/assets/D04.pdf}}
\Box{148}{209}{42}{\sffamily\bfseries\fontsize{8.7}{12}\selectfont\color{gold}\centering Llama-Nemotron}
\Box{22}{220}{166}{\sffamily\bfseries\fontsize{8.8}{13}\selectfont\color{navy}图 1-2　派生、使用与数据接续，共同揭示开放影响}
\Box{22}{227}{166}{\sffamily\fontsize{7.8}{11}\selectfont\color{muted}后训练指基础训练之后的能力调整；样本复用不等于模型参数继承，也不代表性能贡献比例。}
\Box{22}{238}{166}{\sffamily\bfseries\fontsize{9.5}{14}\selectfont\color{teal}开放继续延伸至基础设施}
\Box{22}{245}{166}{\sffamily\fontsize{8.7}{12}\selectfont\color{ink}DeepSeek 开放计算内核、通信、并行训练与存储组件，为模型之外的体系协作提供基础。}

\Note{22}{256}{166}{来源：\href{\HF}{HF Summer 2026}；\href{\ATOM}{ATOM 论文 v1（2026-04）}图 9、12；\href{\NV}{NVIDIA（2025-05-14）}表 2；\href{\DSINFRA}{DeepSeek 官方开放基础设施索引}。}
\Note{22}{265}{166}{口径：Qwen 为 HF 夏季报告观察；47\% 为 ATOM 所追踪 250B+ 模型累计下载份额，按总参数分组；31.1\% 为论文标注“2026 年 1 月”观察点的 OpenRouter 每月前十开放模型词元份额。非全球市场份额，各指标不直接比较；NVIDIA 案例属历史背景，均不并入本报告主期总量。}
\end{canvas}

% Page 4: international foundations and practical public-capability conditions.
\begin{canvas}
\Header{1.2 · 发展目标}{teal}\Footer{04}
\Sec{teal}{1.2　开源发展目标：从开放资源到公共能力}{section12}
\Title{从共同发展，到可持续公共能力}
\Subhead{59}{teal}{明确目标、核对条件，再检验实际结果}
\Body{22}{72}{166}{联合国可持续发展目标（SDGs）提出 17 项共同发展目标。数字公共品（DPG）可包括开源软件、开放数据、人工智能系统与开放内容。数字公共品联盟（DPGA）通过九项要求，评估方案的开放与设计条件；本地实施结果仍需独立检验。}
\Box{22}{101}{78}{\sffamily\bfseries\fontsize{10.6}{15}\selectfont\color{blue}17 项 SDGs｜共同发展方向}
\Box{108}{101}{80}{\sffamily\bfseries\fontsize{10.6}{15}\selectfont\color{teal}9 项 DPG 要求｜开放与设计}
\Rule{22}{111}{101}{blue!25}
\Rule{108}{111}{188}{teal!25}
\definecolor{sdg1}{HTML}{E5243B}
\Rect{22}{115.7}{0.9}{5.7}{sdg1}
\Box{24.5}{115.0}{7}{\sffamily\bfseries\fontsize{7.8}{10}\selectfont\color{muted}01}
\Box{31}{115.0}{30}{\sffamily\fontsize{8.3}{10.8}\selectfont\color{navy}无贫穷}
\definecolor{sdg2}{HTML}{DDA63A}
\Rect{22}{124.5}{0.9}{5.7}{sdg2}
\Box{24.5}{123.8}{7}{\sffamily\bfseries\fontsize{7.8}{10}\selectfont\color{muted}02}
\Box{31}{123.8}{30}{\sffamily\fontsize{8.3}{10.8}\selectfont\color{navy}零饥饿}
\definecolor{sdg3}{HTML}{4C9F38}
\Rect{22}{133.29999999999998}{0.9}{5.7}{sdg3}
\Box{24.5}{132.6}{7}{\sffamily\bfseries\fontsize{7.8}{10}\selectfont\color{muted}03}
\Box{31}{132.6}{30}{\sffamily\fontsize{8.3}{10.8}\selectfont\color{navy}良好健康与福祉}
\definecolor{sdg4}{HTML}{C5192D}
\Rect{22}{142.1}{0.9}{5.7}{sdg4}
\Box{24.5}{141.4}{7}{\sffamily\bfseries\fontsize{7.8}{10}\selectfont\color{muted}04}
\Box{31}{141.4}{30}{\sffamily\fontsize{8.3}{10.8}\selectfont\color{navy}优质教育}
\definecolor{sdg5}{HTML}{FF3A21}
\Rect{22}{150.89999999999998}{0.9}{5.7}{sdg5}
\Box{24.5}{150.2}{7}{\sffamily\bfseries\fontsize{7.8}{10}\selectfont\color{muted}05}
\Box{31}{150.2}{30}{\sffamily\fontsize{8.3}{10.8}\selectfont\color{navy}性别平等}
\definecolor{sdg6}{HTML}{26BDE2}
\Rect{22}{159.7}{0.9}{5.7}{sdg6}
\Box{24.5}{159.0}{7}{\sffamily\bfseries\fontsize{7.8}{10}\selectfont\color{muted}06}
\Box{31}{159.0}{30}{\sffamily\fontsize{8.3}{10.8}\selectfont\color{navy}清洁饮水和卫生设施}
\definecolor{sdg7}{HTML}{FCC30B}
\Rect{22}{168.5}{0.9}{5.7}{sdg7}
\Box{24.5}{167.8}{7}{\sffamily\bfseries\fontsize{7.8}{10}\selectfont\color{muted}07}
\Box{31}{167.8}{30}{\sffamily\fontsize{8.3}{10.8}\selectfont\color{navy}经济适用的清洁能源}
\definecolor{sdg8}{HTML}{A21942}
\Rect{22}{177.3}{0.9}{5.7}{sdg8}
\Box{24.5}{176.60000000000002}{7}{\sffamily\bfseries\fontsize{7.8}{10}\selectfont\color{muted}08}
\Box{31}{176.60000000000002}{30}{\sffamily\fontsize{8.3}{10.8}\selectfont\color{navy}体面工作和经济增长}
\definecolor{sdg9}{HTML}{FD6925}
\Rect{22}{186.1}{0.9}{5.7}{sdg9}
\Box{24.5}{185.4}{7}{\sffamily\bfseries\fontsize{7.8}{10}\selectfont\color{muted}09}
\Box{31}{185.4}{30}{\sffamily\fontsize{8.3}{10.8}\selectfont\color{navy}产业、创新和基础设施}
\definecolor{sdg10}{HTML}{DD1367}
\Rect{63}{115.7}{0.9}{5.7}{sdg10}
\Box{65.5}{115.0}{7}{\sffamily\bfseries\fontsize{7.8}{10}\selectfont\color{muted}10}
\Box{72}{115.0}{30}{\sffamily\fontsize{8.3}{10.8}\selectfont\color{navy}减少不平等}
\definecolor{sdg11}{HTML}{FD9D24}
\Rect{63}{124.5}{0.9}{5.7}{sdg11}
\Box{65.5}{123.8}{7}{\sffamily\bfseries\fontsize{7.8}{10}\selectfont\color{muted}11}
\Box{72}{123.8}{30}{\sffamily\fontsize{8.3}{10.8}\selectfont\color{navy}可持续城市和社区}
\definecolor{sdg12}{HTML}{BF8B2E}
\Rect{63}{133.29999999999998}{0.9}{5.7}{sdg12}
\Box{65.5}{132.6}{7}{\sffamily\bfseries\fontsize{7.8}{10}\selectfont\color{muted}12}
\Box{72}{132.6}{30}{\sffamily\fontsize{8.3}{10.8}\selectfont\color{navy}负责任消费和生产}
\definecolor{sdg13}{HTML}{3F7E44}
\Rect{63}{142.1}{0.9}{5.7}{sdg13}
\Box{65.5}{141.4}{7}{\sffamily\bfseries\fontsize{7.8}{10}\selectfont\color{muted}13}
\Box{72}{141.4}{30}{\sffamily\fontsize{8.3}{10.8}\selectfont\color{navy}气候行动}
\definecolor{sdg14}{HTML}{0A97D9}
\Rect{63}{150.89999999999998}{0.9}{5.7}{sdg14}
\Box{65.5}{150.2}{7}{\sffamily\bfseries\fontsize{7.8}{10}\selectfont\color{muted}14}
\Box{72}{150.2}{30}{\sffamily\fontsize{8.3}{10.8}\selectfont\color{navy}水下生物}
\definecolor{sdg15}{HTML}{56C02B}
\Rect{63}{159.7}{0.9}{5.7}{sdg15}
\Box{65.5}{159.0}{7}{\sffamily\bfseries\fontsize{7.8}{10}\selectfont\color{muted}15}
\Box{72}{159.0}{30}{\sffamily\fontsize{8.3}{10.8}\selectfont\color{navy}陆地生物}
\definecolor{sdg16}{HTML}{00689D}
\Rect{63}{168.5}{0.9}{5.7}{sdg16}
\Box{65.5}{167.8}{7}{\sffamily\bfseries\fontsize{7.8}{10}\selectfont\color{muted}16}
\Box{72}{167.8}{30}{\sffamily\fontsize{8.3}{10.8}\selectfont\color{navy}和平、正义与强大机构}
\definecolor{sdg17}{HTML}{19486A}
\Rect{63}{177.3}{0.9}{5.7}{sdg17}
\Box{65.5}{176.60000000000002}{7}{\sffamily\bfseries\fontsize{7.8}{10}\selectfont\color{muted}17}
\Box{72}{176.60000000000002}{30}{\sffamily\fontsize{8.3}{10.8}\selectfont\color{navy}促进目标实现的伙伴关系}
\Box{108}{115.0}{7}{\sffamily\bfseries\fontsize{8.1}{11}\selectfont\color{teal}01}
\Box{117}{115.0}{32}{\sffamily\bfseries\fontsize{8.6}{11}\selectfont\color{navy}SDG 相关性}
\Box{151}{115.0}{37}{\sffamily\fontsize{7.8}{10.7}\selectfont\color{muted}说明发展目标关联}
\Rule{108}{122.3}{188}{teal!12}
\Box{108}{123.8}{7}{\sffamily\bfseries\fontsize{8.1}{11}\selectfont\color{teal}02}
\Box{117}{123.8}{32}{\sffamily\bfseries\fontsize{8.6}{11}\selectfont\color{navy}开放许可}
\Box{151}{123.8}{37}{\sffamily\fontsize{7.8}{10.7}\selectfont\color{muted}允许合法复用}
\Rule{108}{131.1}{188}{teal!12}
\Box{108}{132.6}{7}{\sffamily\bfseries\fontsize{8.1}{11}\selectfont\color{teal}03}
\Box{117}{132.6}{32}{\sffamily\bfseries\fontsize{8.6}{11}\selectfont\color{navy}清晰权属}
\Box{151}{132.6}{37}{\sffamily\fontsize{7.8}{10.7}\selectfont\color{muted}说明权利归属}
\Rule{108}{139.9}{188}{teal!12}
\Box{108}{141.4}{7}{\sffamily\bfseries\fontsize{8.1}{11}\selectfont\color{teal}04}
\Box{117}{141.4}{32}{\sffamily\bfseries\fontsize{8.6}{11}\selectfont\color{navy}平台独立性}
\Box{151}{141.4}{37}{\sffamily\fontsize{7.8}{10.7}\selectfont\color{muted}减少封闭平台依赖}
\Rule{108}{148.70000000000002}{188}{teal!12}
\Box{108}{150.2}{7}{\sffamily\bfseries\fontsize{8.1}{11}\selectfont\color{teal}05}
\Box{117}{150.2}{32}{\sffamily\bfseries\fontsize{8.6}{11}\selectfont\color{navy}文档}
\Box{151}{150.2}{37}{\sffamily\fontsize{7.8}{10.7}\selectfont\color{muted}支持理解与使用}
\Rule{108}{157.5}{188}{teal!12}
\Box{108}{159.0}{7}{\sffamily\bfseries\fontsize{8.1}{11}\selectfont\color{teal}06}
\Box{117}{159.0}{32}{\sffamily\bfseries\fontsize{8.6}{11}\selectfont\color{navy}数据提取机制}
\Box{151}{159.0}{37}{\sffamily\fontsize{7.8}{10.7}\selectfont\color{muted}可提取非身份识别数据}
\Rule{108}{166.3}{188}{teal!12}
\Box{108}{167.8}{7}{\sffamily\bfseries\fontsize{8.1}{11}\selectfont\color{teal}07}
\Box{117}{167.8}{32}{\sffamily\bfseries\fontsize{8.6}{11}\selectfont\color{navy}隐私与法律}
\Box{151}{167.8}{37}{\sffamily\fontsize{7.8}{10.7}\selectfont\color{muted}遵守隐私及适用法律}
\Rule{108}{175.10000000000002}{188}{teal!12}
\Box{108}{176.60000000000002}{7}{\sffamily\bfseries\fontsize{8.1}{11}\selectfont\color{teal}08}
\Box{117}{176.60000000000002}{32}{\sffamily\bfseries\fontsize{8.6}{11}\selectfont\color{navy}标准与最佳实践}
\Box{151}{176.60000000000002}{37}{\sffamily\fontsize{7.8}{10.7}\selectfont\color{muted}遵循相关标准与实践}
\Rule{108}{183.90000000000003}{188}{teal!12}
\Box{108}{185.4}{7}{\sffamily\bfseries\fontsize{8.1}{11}\selectfont\color{teal}09}
\Box{117}{185.4}{32}{\sffamily\bfseries\fontsize{8.6}{11}\selectfont\color{navy}不造成伤害}
\Box{151}{185.4}{37}{\sffamily\fontsize{7.8}{10.7}\selectfont\color{muted}安全、内容与反骚扰保护}
\Box{22}{199}{166}{\sffamily\bfseries\fontsize{10.5}{15}\selectfont\color{navy}开放资源转化为公共能力，还要跨越四道实施门槛}
\draw[blue!35,line width=0.9pt,-{Stealth[length=1.4mm]}] \Pos{47}{219} .. controls \Pos{58}{216} and \Pos{63}{222} .. \Pos{76}{219};
\fill[blue!5] \Pos{40}{219} circle[radius=6.7mm];
\shade[top color=white,bottom color=blue!28,shading angle=30,draw=blue!35,line width=.45pt] \Pos{40}{219} circle[radius=5.5mm];
\draw[white,line width=.7pt] \Pos{36.81}{215.81} arc[start angle=135,end angle=65,radius=4.51mm];
\Box{37.14}{216.14}{5.720000000000001}{\sffamily\fontsize{9.5}{14}\selectfont\color{ink}\includegraphics[width=5.720000000000001mm]{.build/assets/D05.pdf}}
\Box{22}{227}{36}{\sffamily\bfseries\fontsize{9.8}{14}\selectfont\color{blue}\centering 可获得}
\Box{20}{236}{40}{\sffamily\fontsize{8}{11}\selectfont\color{muted}\centering 版本 · 许可 · 文档}
\draw[teal!35,line width=0.9pt,-{Stealth[length=1.4mm]}] \Pos{90}{219} .. controls \Pos{101}{216} and \Pos{106}{222} .. \Pos{119}{219};
\fill[teal!5] \Pos{83}{219} circle[radius=6.7mm];
\shade[top color=white,bottom color=teal!28,shading angle=30,draw=teal!35,line width=.45pt] \Pos{83}{219} circle[radius=5.5mm];
\draw[white,line width=.7pt] \Pos{79.81}{215.81} arc[start angle=135,end angle=65,radius=4.51mm];
\Box{80.14}{216.14}{5.720000000000001}{\sffamily\fontsize{9.5}{14}\selectfont\color{ink}\includegraphics[width=5.720000000000001mm]{.build/assets/D08.pdf}}
\Box{65}{227}{36}{\sffamily\bfseries\fontsize{9.8}{14}\selectfont\color{teal}\centering 可复用}
\Box{63}{236}{40}{\sffamily\fontsize{8}{11}\selectfont\color{muted}\centering 成本 · 语言 · 工具}
\draw[gold!35,line width=0.9pt,-{Stealth[length=1.4mm]}] \Pos{133}{219} .. controls \Pos{144}{216} and \Pos{149}{222} .. \Pos{162}{219};
\fill[gold!5] \Pos{126}{219} circle[radius=6.7mm];
\shade[top color=white,bottom color=gold!28,shading angle=30,draw=gold!35,line width=.45pt] \Pos{126}{219} circle[radius=5.5mm];
\draw[white,line width=.7pt] \Pos{122.81}{215.81} arc[start angle=135,end angle=65,radius=4.51mm];
\Box{123.14}{216.14}{5.720000000000001}{\sffamily\fontsize{9.5}{14}\selectfont\color{ink}\includegraphics[width=5.720000000000001mm]{.build/assets/D09.pdf}}
\Box{108}{227}{36}{\sffamily\bfseries\fontsize{9.8}{14}\selectfont\color{gold}\centering 可维护}
\Box{106}{236}{40}{\sffamily\fontsize{8}{11}\selectfont\color{muted}\centering 责任 · 资金 · 安全}
\fill[violet!5] \Pos{169}{219} circle[radius=6.7mm];
\shade[top color=white,bottom color=violet!28,shading angle=30,draw=violet!35,line width=.45pt] \Pos{169}{219} circle[radius=5.5mm];
\draw[white,line width=.7pt] \Pos{165.81}{215.81} arc[start angle=135,end angle=65,radius=4.51mm];
\Box{166.14}{216.14}{5.720000000000001}{\sffamily\fontsize{9.5}{14}\selectfont\color{ink}\includegraphics[width=5.720000000000001mm]{.build/assets/D10.pdf}}
\Box{151}{227}{36}{\sffamily\bfseries\fontsize{9.8}{14}\selectfont\color{violet}\centering 有成效}
\Box{149}{236}{40}{\sffamily\fontsize{8}{11}\selectfont\color{muted}\centering 覆盖 · 质量 · 公平}

\Body{22}{247}{166}{AI 普惠需要本地伙伴参与问题定义与技术改进，也需要基金会、企业、研究机构和社区共同承担长期维护。资源开放之后，谁能够使用、谁负责维护、谁持续受益，应成为公共投入的共同问题。}
\Note{22}{268}{166}{图 1-3　来源：\href{\UNGoals}{联合国 SDGs}、\href{\DPG}{DPGA 标准}（2026-11-09 核对；中文简述与改绘部分 CC BY-SA 4.0）。两套体系不逐项对应；四道门槛为本报告分析框架，非自动生效的因果链。}
\end{canvas}

% Page 5: approved signature and complete original framework.
\begin{canvas}
\Header{1.2 · 发展目标}{teal}\Footer{05}
\Kicker{teal}{1.2 / OPEN SOURCE DEVELOPMENT GOALS}
\Title{OSDGs：开源发展的共同语言}
\Subhead{59}{teal}{让人才、供给、连接与治理共同支撑公共价值}
\Body{22}{72}{166}{本报告提出开源发展目标（OSDGs），将全球开源的增长与结构变化放回发展问题：谁能参与，什么可以复用，成果如何接续，合作怎样持续。四个支柱互为补充，共同构成观察人工智能时代开放创新的框架。}
\draw[blue!12,line width=5pt] \Pos{79}{125} .. controls \Pos{88}{125} and \Pos{84}{135} .. \Pos{91}{142};
\draw[blue!55,line width=0.7pt] \Pos{79}{125} .. controls \Pos{88}{125} and \Pos{84}{135} .. \Pos{91}{142};
\draw[teal!12,line width=5pt] \Pos{131}{125} .. controls \Pos{122}{125} and \Pos{127}{135} .. \Pos{119}{142};
\draw[teal!55,line width=0.7pt] \Pos{131}{125} .. controls \Pos{122}{125} and \Pos{127}{135} .. \Pos{119}{142};
\draw[gold!12,line width=5pt] \Pos{79}{184} .. controls \Pos{88}{184} and \Pos{84}{172} .. \Pos{91}{165};
\draw[gold!55,line width=0.7pt] \Pos{79}{184} .. controls \Pos{88}{184} and \Pos{84}{172} .. \Pos{91}{165};
\draw[violet!12,line width=5pt] \Pos{131}{184} .. controls \Pos{122}{184} and \Pos{127}{172} .. \Pos{119}{165};
\draw[violet!55,line width=0.7pt] \Pos{131}{184} .. controls \Pos{122}{184} and \Pos{127}{172} .. \Pos{119}{165};
\Box{86}{135}{38}{\sffamily\fontsize{9.5}{14}\selectfont\color{ink}\includegraphics[width=38mm]{.build/assets/osdgs-display.pdf}}
\Box{79}{177}{52}{\sffamily\bfseries\fontsize{9.1}{13}\selectfont\color{teal}\centering 共同发展 · 全球接续}
\shade[left color=blue!10,right color=blue!2,rounded corners=2mm,draw=blue!17,line width=.45pt] \Pos{22}{99} rectangle \Pos{80}{152};
\path[fill=blue!9] \Pos{42.879999999999995}{152} .. controls \Pos{61.440000000000005}{151} and \Pos{66.08}{147} .. \Pos{79}{148} -- \Pos{79}{151} -- \Pos{42.879999999999995}{152} -- cycle;
\Box{25}{102}{52}{\sffamily\bfseries\fontsize{9.6}{14}\selectfont\color{blue}X1 人才规模与多样性}
\Box{25}{111}{52}{\sffamily\fontsize{8.4}{12}\selectfont\color{muted}扩大参与与持续成长}
\Box{25}{123.0}{4.2}{\sffamily\fontsize{9.5}{14}\selectfont\color{ink}\includegraphics[width=4.2mm]{.build/assets/D01.pdf}}
\Box{32}{123.0}{45}{\sffamily\fontsize{8.6}{12}\selectfont\color{ink}D01 规模与密度}
\Box{25}{131.7}{4.2}{\sffamily\fontsize{9.5}{14}\selectfont\color{ink}\includegraphics[width=4.2mm]{.build/assets/D02.pdf}}
\Box{32}{131.7}{45}{\sffamily\fontsize{8.6}{12}\selectfont\color{ink}D02 生态多样性}
\shade[left color=teal!10,right color=teal!2,rounded corners=2mm,draw=teal!17,line width=.45pt] \Pos{130}{99} rectangle \Pos{188}{152};
\path[fill=teal!9] \Pos{150.88}{152} .. controls \Pos{169.44}{151} and \Pos{174.07999999999998}{147} .. \Pos{187}{148} -- \Pos{187}{151} -- \Pos{150.88}{152} -- cycle;
\Box{133}{102}{52}{\sffamily\bfseries\fontsize{9.6}{14}\selectfont\color{teal}X2 技术产出与开放供给}
\Box{133}{111}{52}{\sffamily\fontsize{8.4}{12}\selectfont\color{muted}形成可复用的开放资源}
\Box{133}{123.0}{4.2}{\sffamily\fontsize{9.5}{14}\selectfont\color{ink}\includegraphics[width=4.2mm]{.build/assets/D03.pdf}}
\Box{140}{123.0}{45}{\sffamily\fontsize{8.6}{12}\selectfont\color{ink}D03 代码产出}
\Box{133}{131.7}{4.2}{\sffamily\fontsize{9.5}{14}\selectfont\color{ink}\includegraphics[width=4.2mm]{.build/assets/D04.pdf}}
\Box{140}{131.7}{45}{\sffamily\fontsize{8.6}{12}\selectfont\color{ink}D04 模型开放供给}
\Box{133}{140.4}{4.2}{\sffamily\fontsize{9.5}{14}\selectfont\color{ink}\includegraphics[width=4.2mm]{.build/assets/D05.pdf}}
\Box{140}{140.4}{45}{\sffamily\fontsize{8.6}{12}\selectfont\color{ink}D05 数据开放供给}
\shade[left color=gold!10,right color=gold!2,rounded corners=2mm,draw=gold!17,line width=.45pt] \Pos{22}{160} rectangle \Pos{80}{213};
\path[fill=gold!9] \Pos{42.879999999999995}{213} .. controls \Pos{61.440000000000005}{212} and \Pos{66.08}{208} .. \Pos{79}{209} -- \Pos{79}{212} -- \Pos{42.879999999999995}{213} -- cycle;
\Box{25}{163}{52}{\sffamily\bfseries\fontsize{9.6}{14}\selectfont\color{gold}X3 全球连接与网络效应}
\Box{25}{172}{52}{\sffamily\fontsize{8.4}{12}\selectfont\color{muted}让成果跨境接续生长}
\Box{25}{184.0}{4.2}{\sffamily\fontsize{9.5}{14}\selectfont\color{ink}\includegraphics[width=4.2mm]{.build/assets/D06.pdf}}
\Box{32}{184.0}{45}{\sffamily\fontsize{8.6}{12}\selectfont\color{ink}D06 代码影响力}
\Box{25}{192.7}{4.2}{\sffamily\fontsize{9.5}{14}\selectfont\color{ink}\includegraphics[width=4.2mm]{.build/assets/D07.pdf}}
\Box{32}{192.7}{45}{\sffamily\fontsize{8.6}{12}\selectfont\color{ink}D07 模型影响力}
\Box{25}{201.4}{4.2}{\sffamily\fontsize{9.5}{14}\selectfont\color{ink}\includegraphics[width=4.2mm]{.build/assets/D08.pdf}}
\Box{32}{201.4}{45}{\sffamily\fontsize{8.6}{12}\selectfont\color{ink}D08 数据集影响力}
\shade[left color=violet!10,right color=violet!2,rounded corners=2mm,draw=violet!17,line width=.45pt] \Pos{130}{160} rectangle \Pos{188}{213};
\path[fill=violet!9] \Pos{150.88}{213} .. controls \Pos{169.44}{212} and \Pos{174.07999999999998}{208} .. \Pos{187}{209} -- \Pos{187}{212} -- \Pos{150.88}{213} -- cycle;
\Box{133}{163}{52}{\sffamily\bfseries\fontsize{9.6}{14}\selectfont\color{violet}X4 开源善治与公共价值}
\Box{133}{172}{52}{\sffamily\fontsize{8.4}{12}\selectfont\color{muted}让协作持续、成果共享}
\Box{133}{184.0}{4.2}{\sffamily\fontsize{9.5}{14}\selectfont\color{ink}\includegraphics[width=4.2mm]{.build/assets/D09.pdf}}
\Box{140}{184.0}{45}{\sffamily\fontsize{8.6}{12}\selectfont\color{ink}D09 治理成熟度}
\Box{133}{192.7}{4.2}{\sffamily\fontsize{9.5}{14}\selectfont\color{ink}\includegraphics[width=4.2mm]{.build/assets/D10.pdf}}
\Box{140}{192.7}{45}{\sffamily\fontsize{8.6}{12}\selectfont\color{ink}D10 公共福祉}
\Box{22}{218}{166}{\sffamily\fontsize{8.2}{12}\selectfont\color{muted}\centering OSDGs / Open Source Development Goals}

\Cap{227}{图 1-4　开源发展目标：四大战略支柱与十个观察维度}
\Body{22}{237}{166}{\textbf{人才与供给相互支撑。} 参与机会需要转化为持续贡献，开放成果需要具备可理解、可使用的条件。}
\Body{22}{251}{166}{\textbf{连接与治理共同延长价值。} 把跨境接续、共同维护与受益群体纳入观察，才能解释资源增长之后的体系能力。}
\Note{22}{270}{166}{来源：本报告获批框架与动态星座 v2.0.0。OSDGs 是本报告研究框架，不是联合国目标或认证；十维编号用于导视，不是测算权重。}
\end{canvas}

% Page 6: original measurement system and relational observation in one view.
\begin{canvas}
\Header{1.3 · 观察方法}{gold}\Footer{06}
\Sec{gold}{1.3　观察方法：跨平台的全球开源全景}{section13}
\Title{OSDI：从发展目标到关系证据}
\Subhead{59}{gold}{让目标可观察，让分散的对象彼此联系}
\Body{22}{72}{166}{本报告构建开源发展指标体系（OSDI），将 OSDGs 分解为观察方向和候选指标。例如，人才目标可观察活跃开发者、新人进入与机构贡献。科学性要求指标对应问题，可操作性要求定义可核，前瞻性要求覆盖模型、数据与复用。}
\shade[left color=blue!10,right color=blue!2,rounded corners=2mm,draw=blue!17,line width=.45pt] \Pos{22}{99} rectangle \Pos{188}{158};
\path[fill=blue!9] \Pos{81.75999999999999}{158} .. controls \Pos{134.88}{157} and \Pos{148.16}{153} .. \Pos{187}{154} -- \Pos{187}{157} -- \Pos{81.75999999999999}{158} -- cycle;
\Box{27}{102}{35}{\sffamily\bfseries\fontsize{18}{23}\selectfont\color{navy}OSDI}
\Box{63}{104}{120}{\sffamily\bfseries\fontsize{11}{15}\selectfont\color{blue}开源发展指标体系｜把目标转化为观察信号}
\Box{33}{113}{34}{\sffamily\bfseries\fontsize{27}{31}\selectfont\color{blue}\centering 4}
\Box{30}{133}{40}{\sffamily\bfseries\fontsize{10}{14}\selectfont\color{navy}\centering 发展目标}
\LinkLine{73}{126}{82}{126}{blue!30}
\Box{88}{113}{34}{\sffamily\bfseries\fontsize{27}{31}\selectfont\color{teal}\centering 10}
\Box{85}{133}{40}{\sffamily\bfseries\fontsize{10}{14}\selectfont\color{navy}\centering 观察方向}
\LinkLine{128}{126}{137}{126}{teal!30}
\Box{143}{113}{34}{\sffamily\bfseries\fontsize{27}{31}\selectfont\color{violet}\centering 29}
\Box{140}{133}{40}{\sffamily\bfseries\fontsize{10}{14}\selectfont\color{navy}\centering 候选指标}
\Box{27}{147}{156}{\sffamily\fontsize{8.4}{12}\selectfont\color{muted}X1 人才：5 项　　X2 供给：7 项　　X3 连接：11 项　　X4 治理与公共价值：6 项}
\Box{22}{163}{50}{\sffamily\bfseries\fontsize{10}{14}\selectfont\color{blue}北极星}
\Box{40}{163}{33}{\sffamily\fontsize{8.5}{12}\selectfont\color{muted}把握长期方向}
\Box{80}{163}{50}{\sffamily\bfseries\fontsize{10}{14}\selectfont\color{teal}诊断}
\Box{98}{163}{33}{\sffamily\fontsize{8.5}{12}\selectfont\color{muted}识别差异与瓶颈}
\Box{138}{163}{50}{\sffamily\bfseries\fontsize{10}{14}\selectfont\color{gold}驱动}
\Box{156}{163}{33}{\sffamily\fontsize{8.5}{12}\selectfont\color{muted}跟踪过程与变化}

\Note{22}{175}{166}{图 1-5　OSDI 的结构与用途。29 项为候选目录；“北极星、诊断、驱动”是用途，不是证据等级。}
\Box{22}{188}{166}{\sffamily\bfseries\fontsize{10.5}{15}\selectfont\color{navy}把指标放回协作关系：六类对象联读}
\draw[blue!45,line width=0.7pt,-{Stealth[length=1.4mm]}] \Pos{49}{204} .. controls \Pos{65}{201} and \Pos{81}{207} .. \Pos{97}{204};
\draw[blue!45,line width=0.7pt,-{Stealth[length=1.4mm]}] \Pos{108}{204} .. controls \Pos{124}{201} and \Pos{140}{207} .. \Pos{156}{204};
\draw[teal!45,line width=0.7pt,-{Stealth[length=1.4mm]}] \Pos{161}{210} .. controls \Pos{161}{229} and \Pos{48}{219} .. \Pos{48}{235};
\draw[gold!45,line width=0.7pt,-{Stealth[length=1.4mm]}] \Pos{161}{210} .. controls \Pos{161}{222} and \Pos{103}{222} .. \Pos{103}{235};
\draw[violet!45,line width=0.7pt,-{Stealth[length=1.4mm]}] \Pos{161}{210} .. controls \Pos{161}{220} and \Pos{161}{226} .. \Pos{161}{235};
\fill[blue!5] \Pos{42}{204} circle[radius=6.4mm];
\shade[top color=white,bottom color=blue!28,shading angle=30,draw=blue!35,line width=.45pt] \Pos{42}{204} circle[radius=5.2mm];
\draw[white,line width=.7pt] \Pos{38.984}{200.984} arc[start angle=135,end angle=65,radius=4.264mm];
\Box{39.296}{201.296}{5.408}{\sffamily\fontsize{9.5}{14}\selectfont\color{ink}\includegraphics[width=5.408mm]{.build/assets/D01.pdf}}
\Box{22.0}{211}{40}{\sffamily\bfseries\fontsize{10}{14}\selectfont\color{blue}\centering 人才}
\fill[blue!5] \Pos{103}{204} circle[radius=6.4mm];
\shade[top color=white,bottom color=blue!28,shading angle=30,draw=blue!35,line width=.45pt] \Pos{103}{204} circle[radius=5.2mm];
\draw[white,line width=.7pt] \Pos{99.984}{200.984} arc[start angle=135,end angle=65,radius=4.264mm];
\Box{100.296}{201.296}{5.408}{\sffamily\fontsize{9.5}{14}\selectfont\color{ink}\includegraphics[width=5.408mm]{.build/assets/D02.pdf}}
\Box{83.0}{211}{40}{\sffamily\bfseries\fontsize{10}{14}\selectfont\color{blue}\centering 组织}
\fill[blue!5] \Pos{162}{204} circle[radius=6.4mm];
\shade[top color=white,bottom color=blue!28,shading angle=30,draw=blue!35,line width=.45pt] \Pos{162}{204} circle[radius=5.2mm];
\draw[white,line width=.7pt] \Pos{158.984}{200.984} arc[start angle=135,end angle=65,radius=4.264mm];
\Box{159.296}{201.296}{5.408}{\sffamily\fontsize{9.5}{14}\selectfont\color{ink}\includegraphics[width=5.408mm]{.build/assets/D03.pdf}}
\Box{142.0}{211}{40}{\sffamily\bfseries\fontsize{10}{14}\selectfont\color{blue}\centering 项目}
\fill[teal!5] \Pos{42}{239} circle[radius=6.4mm];
\shade[top color=white,bottom color=teal!28,shading angle=30,draw=teal!35,line width=.45pt] \Pos{42}{239} circle[radius=5.2mm];
\draw[white,line width=.7pt] \Pos{38.984}{235.984} arc[start angle=135,end angle=65,radius=4.264mm];
\Box{39.296}{236.296}{5.408}{\sffamily\fontsize{9.5}{14}\selectfont\color{ink}\includegraphics[width=5.408mm]{.build/assets/D03.pdf}}
\Box{22.0}{246}{40}{\sffamily\bfseries\fontsize{10}{14}\selectfont\color{teal}\centering 代码}
\fill[gold!5] \Pos{103}{239} circle[radius=6.4mm];
\shade[top color=white,bottom color=gold!28,shading angle=30,draw=gold!35,line width=.45pt] \Pos{103}{239} circle[radius=5.2mm];
\draw[white,line width=.7pt] \Pos{99.984}{235.984} arc[start angle=135,end angle=65,radius=4.264mm];
\Box{100.296}{236.296}{5.408}{\sffamily\fontsize{9.5}{14}\selectfont\color{ink}\includegraphics[width=5.408mm]{.build/assets/D04.pdf}}
\Box{83.0}{246}{40}{\sffamily\bfseries\fontsize{10}{14}\selectfont\color{gold}\centering 模型}
\fill[violet!5] \Pos{162}{239} circle[radius=6.4mm];
\shade[top color=white,bottom color=violet!28,shading angle=30,draw=violet!35,line width=.45pt] \Pos{162}{239} circle[radius=5.2mm];
\draw[white,line width=.7pt] \Pos{158.984}{235.984} arc[start angle=135,end angle=65,radius=4.264mm];
\Box{159.296}{236.296}{5.408}{\sffamily\fontsize{9.5}{14}\selectfont\color{ink}\includegraphics[width=5.408mm]{.build/assets/D05.pdf}}
\Box{142.0}{246}{40}{\sffamily\bfseries\fontsize{10}{14}\selectfont\color{violet}\centering 数据}
\Box{64}{197}{18}{\sffamily\fontsize{7.5}{10}\selectfont\color{muted}\centering 参与}
\Box{125}{197}{18}{\sffamily\fontsize{7.5}{10}\selectfont\color{muted}\centering 共建}
\Box{78}{225}{61}{\sffamily\fontsize{8.1}{11}\selectfont\color{muted}\centering 项目形成与维护开放资产}

\Note{22}{260}{166}{图 1-6　六类对象的关联观察示意。来源：2026 版研究框架与\href{\OD}{OpenDigger} 方法材料。身份与关系仅在已核验范围内联读，不表示完整测量；各类资产保留独立统计单位，完整指标定义见附录 A。}
\end{canvas}

% Page 7: accessible explanation, explicit sample/period boundaries.
\begin{canvas}
\Header{1.3 · 观察方法}{gold}\Footer{07}
\Kicker{gold}{1.3 / PARTICIPATION, CONTRIBUTION AND INFLUENCE}
\Title{同一生态，需要不同的观察尺度}
\Subhead{59}{gold}{参与、贡献与影响力，共同解释协作}
\Body{22}{72}{166}{参与说明谁进入共同工作，贡献关注完成了哪些协作，影响力观察这些工作在关系网络中的位置。AI 生成与自动化进一步改变了活动量的含义，因此需要辨别参与主体、协作类型与持续性，不能由提交数量直接判断能力。}
\Box{22}{102}{166}{\sffamily\bfseries\fontsize{10.4}{15}\selectfont\color{navy}同一案例：项目 P 的兼容性修复，在项目 Q 中继续使用}
\draw[blue!35,line width=1pt,-{Stealth[length=1.4mm]}] \Pos{48}{130} .. controls \Pos{56}{126} and \Pos{65}{134} .. \Pos{75}{130};
\fill[blue!5] \Pos{40}{130} circle[radius=7.7mm];
\shade[top color=white,bottom color=blue!28,shading angle=30,draw=blue!35,line width=.45pt] \Pos{40}{130} circle[radius=6.5mm];
\draw[white,line width=.7pt] \Pos{36.230000000000004}{126.23} arc[start angle=135,end angle=65,radius=5.33mm];
\Box{36.62}{126.62}{6.76}{\sffamily\fontsize{9.5}{14}\selectfont\color{ink}\includegraphics[width=6.76mm]{.build/assets/D01.pdf}}
\Box{21}{140}{38}{\sffamily\bfseries\fontsize{9.2}{13}\selectfont\color{blue}\centering 甲｜发现问题}
\Box{20}{149}{40}{\sffamily\fontsize{8}{11}\selectfont\color{muted}\centering 报告不兼容现象}
\draw[teal!35,line width=1pt,-{Stealth[length=1.4mm]}] \Pos{91}{130} .. controls \Pos{99}{126} and \Pos{108}{134} .. \Pos{118}{130};
\fill[teal!5] \Pos{83}{130} circle[radius=7.7mm];
\shade[top color=white,bottom color=teal!28,shading angle=30,draw=teal!35,line width=.45pt] \Pos{83}{130} circle[radius=6.5mm];
\draw[white,line width=.7pt] \Pos{79.23}{126.23} arc[start angle=135,end angle=65,radius=5.33mm];
\Box{79.62}{126.62}{6.76}{\sffamily\fontsize{9.5}{14}\selectfont\color{ink}\includegraphics[width=6.76mm]{.build/assets/D03.pdf}}
\Box{64}{140}{38}{\sffamily\bfseries\fontsize{9.2}{13}\selectfont\color{teal}\centering 乙｜提交修复}
\Box{63}{149}{40}{\sffamily\fontsize{8}{11}\selectfont\color{muted}\centering 贡献代码改动}
\draw[gold!35,line width=1pt,-{Stealth[length=1.4mm]}] \Pos{134}{130} .. controls \Pos{142}{126} and \Pos{151}{134} .. \Pos{161}{130};
\fill[gold!5] \Pos{126}{130} circle[radius=7.7mm];
\shade[top color=white,bottom color=gold!28,shading angle=30,draw=gold!35,line width=.45pt] \Pos{126}{130} circle[radius=6.5mm];
\draw[white,line width=.7pt] \Pos{122.23}{126.23} arc[start angle=135,end angle=65,radius=5.33mm];
\Box{122.62}{126.62}{6.76}{\sffamily\fontsize{9.5}{14}\selectfont\color{ink}\includegraphics[width=6.76mm]{.build/assets/D09.pdf}}
\Box{107}{140}{38}{\sffamily\bfseries\fontsize{9.2}{13}\selectfont\color{gold}\centering 丙｜审阅验证}
\Box{106}{149}{40}{\sffamily\fontsize{8}{11}\selectfont\color{muted}\centering 测试并完善修改}
\fill[violet!5] \Pos{169}{130} circle[radius=7.7mm];
\shade[top color=white,bottom color=violet!28,shading angle=30,draw=violet!35,line width=.45pt] \Pos{169}{130} circle[radius=6.5mm];
\draw[white,line width=.7pt] \Pos{165.23}{126.23} arc[start angle=135,end angle=65,radius=5.33mm];
\Box{165.62}{126.62}{6.76}{\sffamily\fontsize{9.5}{14}\selectfont\color{ink}\includegraphics[width=6.76mm]{.build/assets/D08.pdf}}
\Box{150}{140}{38}{\sffamily\bfseries\fontsize{9.2}{13}\selectfont\color{violet}\centering 丁｜接续使用}
\Box{149}{149}{40}{\sffamily\fontsize{8}{11}\selectfont\color{muted}\centering 在项目 Q 中复用}
\shade[left color=blue!10,right color=blue!2,rounded corners=2mm,draw=blue!17,line width=.45pt] \Pos{22}{172} rectangle \Pos{74}{215};
\path[fill=blue!9] \Pos{40.72}{215} .. controls \Pos{57.36}{214} and \Pos{61.52}{210} .. \Pos{73}{211} -- \Pos{73}{214} -- \Pos{40.72}{215} -- cycle;
\fill[blue!5] \Pos{30}{181} circle[radius=5.2mm];
\shade[top color=white,bottom color=blue!28,shading angle=30,draw=blue!35,line width=.45pt] \Pos{30}{181} circle[radius=4mm];
\draw[white,line width=.7pt] \Pos{27.68}{178.68} arc[start angle=135,end angle=65,radius=3.28mm];
\Box{27.92}{178.92}{4.16}{\sffamily\fontsize{9.5}{14}\selectfont\color{ink}\includegraphics[width=4.16mm]{.build/assets/D01.pdf}}
\Box{38}{176}{32}{\sffamily\bfseries\fontsize{11}{15}\selectfont\color{blue}参与}
\Box{26}{190}{44}{\sffamily\fontsize{8.8}{12}\selectfont\color{navy}谁持续进入协作？}
\Box{26}{202}{44}{\sffamily\fontsize{8.3}{12}\selectfont\color{muted}人数、留存与角色}
\shade[left color=teal!10,right color=teal!2,rounded corners=2mm,draw=teal!17,line width=.45pt] \Pos{79}{172} rectangle \Pos{131}{215};
\path[fill=teal!9] \Pos{97.72}{215} .. controls \Pos{114.36}{214} and \Pos{118.52000000000001}{210} .. \Pos{130}{211} -- \Pos{130}{214} -- \Pos{97.72}{215} -- cycle;
\fill[teal!5] \Pos{87}{181} circle[radius=5.2mm];
\shade[top color=white,bottom color=teal!28,shading angle=30,draw=teal!35,line width=.45pt] \Pos{87}{181} circle[radius=4mm];
\draw[white,line width=.7pt] \Pos{84.68}{178.68} arc[start angle=135,end angle=65,radius=3.28mm];
\Box{84.92}{178.92}{4.16}{\sffamily\fontsize{9.5}{14}\selectfont\color{ink}\includegraphics[width=4.16mm]{.build/assets/D03.pdf}}
\Box{95}{176}{32}{\sffamily\bfseries\fontsize{11}{15}\selectfont\color{teal}贡献}
\Box{83}{190}{44}{\sffamily\fontsize{8.8}{12}\selectfont\color{navy}完成了哪些共同工作？}
\Box{83}{202}{44}{\sffamily\fontsize{8.3}{12}\selectfont\color{muted}开发、讨论与审阅}
\shade[left color=gold!10,right color=gold!2,rounded corners=2mm,draw=gold!17,line width=.45pt] \Pos{136}{172} rectangle \Pos{188}{215};
\path[fill=gold!9] \Pos{154.72}{215} .. controls \Pos{171.36}{214} and \Pos{175.52}{210} .. \Pos{187}{211} -- \Pos{187}{214} -- \Pos{154.72}{215} -- cycle;
\fill[gold!5] \Pos{144}{181} circle[radius=5.2mm];
\shade[top color=white,bottom color=gold!28,shading angle=30,draw=gold!35,line width=.45pt] \Pos{144}{181} circle[radius=4mm];
\draw[white,line width=.7pt] \Pos{141.68}{178.68} arc[start angle=135,end angle=65,radius=3.28mm];
\Box{141.92}{178.92}{4.16}{\sffamily\fontsize{9.5}{14}\selectfont\color{ink}\includegraphics[width=4.16mm]{.build/assets/D06.pdf}}
\Box{152}{176}{32}{\sffamily\bfseries\fontsize{11}{15}\selectfont\color{gold}影响力}
\Box{140}{190}{44}{\sffamily\fontsize{8.8}{12}\selectfont\color{navy}工作怎样被连接与接续？}
\Box{140}{202}{44}{\sffamily\fontsize{8.3}{12}\selectfont\color{muted}协作关系与结构位置}
\Box{22}{219}{166}{\sffamily\fontsize{8.1}{12}\selectfont\color{muted}图 1-7　协作示意，非实测样本。OpenRank 综合协作关系，观察贡献在网络中的结构性影响。}

\Rule{22}{230}{188}{line}
\Label{22}{236}{166}{主统计期：2025-07-01 至 2026-06-30}
\Note{22}{247}{166}{历史比较使用可比窗口；外部平台年度窗和期后观察分别标注。跨平台身份、自动化账号、地区未知与项目边界按实际核验状态处理，未知不按零计算。}
\Note{22}{262}{166}{来源：2026 版第一章与\href{\OD}{OpenDigger／OpenRank 方法材料}。OSDI 将四目标、十方向落实为候选指标；完整定义与可得性在附录 A 及研究记录中维护。下载、关注和派生也须分别解释，不能直接等同公共结果。}
\end{canvas}

% Page 8: four complementary perspectives and research questions, no premature findings.
\begin{canvas}
\Header{1.3 · 全书路径}{gold}\Footer{08}
\Kicker{gold}{1.3 / FROM QUESTIONS TO EVIDENCE AND ACTION}
\Title{从中国出发，理解全球开放协作}
\Subhead{59}{gold}{四类证据共同回答人工智能时代的发展问题}
\Body{22}{72}{166}{中国的模型供给、产业与人才实践，为观察全球开源提供了重要入口。本报告将其放入跨境复用与共同创新中，也关注其他地区的长期积累和全球南方的参与路径。地域特色与时序变化，需要由互补证据共同解释。}
\draw[blue!20,line width=3pt] \Pos{60}{123} .. controls \Pos{60}{147} and \Pos{87}{147} .. \Pos{105}{147};
\draw[gold!20,line width=3pt] \Pos{151}{123} .. controls \Pos{151}{147} and \Pos{123}{147} .. \Pos{105}{147};
\draw[teal!20,line width=3pt] \Pos{60}{170} .. controls \Pos{60}{147} and \Pos{87}{147} .. \Pos{105}{147};
\draw[violet!20,line width=3pt] \Pos{151}{170} .. controls \Pos{151}{147} and \Pos{123}{147} .. \Pos{105}{147};
\fill[blue!5] \Pos{105}{147} circle[radius=14.2mm];
\shade[top color=white,bottom color=blue!28,shading angle=30,draw=blue!35,line width=.45pt] \Pos{105}{147} circle[radius=13mm];
\draw[white,line width=.7pt] \Pos{97.46000000000001}{139.46} arc[start angle=135,end angle=65,radius=10.66mm];
\Box{93}{138}{24}{\sffamily\bfseries\fontsize{10}{14}\selectfont\color{navy}\centering 全球开放}
\Box{93}{149}{24}{\sffamily\bfseries\fontsize{12}{16}\selectfont\color{teal}\centering 协作}
\shade[left color=blue!10,right color=blue!2,rounded corners=2mm,draw=blue!17,line width=.45pt] \Pos{22}{103} rectangle \Pos{78}{139};
\path[fill=blue!9] \Pos{42.16}{139} .. controls \Pos{60.080000000000005}{138} and \Pos{64.56}{134} .. \Pos{77}{135} -- \Pos{77}{138} -- \Pos{42.16}{139} -- cycle;
\fill[blue!5] \Pos{31}{113} circle[radius=5.7mm];
\shade[top color=white,bottom color=blue!28,shading angle=30,draw=blue!35,line width=.45pt] \Pos{31}{113} circle[radius=4.5mm];
\draw[white,line width=.7pt] \Pos{28.39}{110.39} arc[start angle=135,end angle=65,radius=3.69mm];
\Box{28.66}{110.66}{4.68}{\sffamily\fontsize{9.5}{14}\selectfont\color{ink}\includegraphics[width=4.68mm]{.build/assets/D01.pdf}}
\Box{39}{107}{37}{\sffamily\bfseries\fontsize{10}{14}\selectfont\color{blue}02 生态地图}
\Box{39}{117}{37}{\sffamily\fontsize{8}{11}\selectfont\color{muted}GOSO}
\Box{26}{129}{48}{\sffamily\fontsize{8.5}{12}\selectfont\color{ink}\centering 参与与资源分布}
\shade[left color=gold!10,right color=gold!2,rounded corners=2mm,draw=gold!17,line width=.45pt] \Pos{132}{103} rectangle \Pos{188}{139};
\path[fill=gold!9] \Pos{152.16}{139} .. controls \Pos{170.08}{138} and \Pos{174.56}{134} .. \Pos{187}{135} -- \Pos{187}{138} -- \Pos{152.16}{139} -- cycle;
\fill[gold!5] \Pos{141}{113} circle[radius=5.7mm];
\shade[top color=white,bottom color=gold!28,shading angle=30,draw=gold!35,line width=.45pt] \Pos{141}{113} circle[radius=4.5mm];
\draw[white,line width=.7pt] \Pos{138.39}{110.39} arc[start angle=135,end angle=65,radius=3.69mm];
\Box{138.66}{110.66}{4.68}{\sffamily\fontsize{9.5}{14}\selectfont\color{ink}\includegraphics[width=4.68mm]{.build/assets/D06.pdf}}
\Box{149}{107}{37}{\sffamily\bfseries\fontsize{10}{14}\selectfont\color{gold}03 影响力排行}
\Box{149}{117}{37}{\sffamily\fontsize{8}{11}\selectfont\color{muted}GOSR}
\Box{136}{129}{48}{\sffamily\fontsize{8.5}{12}\selectfont\color{ink}\centering 关键主体与位置}
\shade[left color=teal!10,right color=teal!2,rounded corners=2mm,draw=teal!17,line width=.45pt] \Pos{22}{155} rectangle \Pos{78}{191};
\path[fill=teal!9] \Pos{42.16}{191} .. controls \Pos{60.080000000000005}{190} and \Pos{64.56}{186} .. \Pos{77}{187} -- \Pos{77}{190} -- \Pos{42.16}{191} -- cycle;
\fill[teal!5] \Pos{31}{165} circle[radius=5.7mm];
\shade[top color=white,bottom color=teal!28,shading angle=30,draw=teal!35,line width=.45pt] \Pos{31}{165} circle[radius=4.5mm];
\draw[white,line width=.7pt] \Pos{28.39}{162.39} arc[start angle=135,end angle=65,radius=3.69mm];
\Box{28.66}{162.66}{4.68}{\sffamily\fontsize{9.5}{14}\selectfont\color{ink}\includegraphics[width=4.68mm]{.build/assets/D03.pdf}}
\Box{39}{159}{37}{\sffamily\bfseries\fontsize{10}{14}\selectfont\color{teal}04 技术全景}
\Box{39}{169}{37}{\sffamily\fontsize{8}{11}\selectfont\color{muted}GOSL}
\Box{26}{181}{48}{\sffamily\fontsize{8.5}{12}\selectfont\color{ink}\centering 技术分层与连接}
\shade[left color=violet!10,right color=violet!2,rounded corners=2mm,draw=violet!17,line width=.45pt] \Pos{132}{155} rectangle \Pos{188}{191};
\path[fill=violet!9] \Pos{152.16}{191} .. controls \Pos{170.08}{190} and \Pos{174.56}{186} .. \Pos{187}{187} -- \Pos{187}{190} -- \Pos{152.16}{191} -- cycle;
\fill[violet!5] \Pos{141}{165} circle[radius=5.7mm];
\shade[top color=white,bottom color=violet!28,shading angle=30,draw=violet!35,line width=.45pt] \Pos{141}{165} circle[radius=4.5mm];
\draw[white,line width=.7pt] \Pos{138.39}{162.39} arc[start angle=135,end angle=65,radius=3.69mm];
\Box{138.66}{162.66}{4.68}{\sffamily\fontsize{9.5}{14}\selectfont\color{ink}\includegraphics[width=4.68mm]{.build/assets/D04.pdf}}
\Box{149}{159}{37}{\sffamily\bfseries\fontsize{10}{14}\selectfont\color{violet}05 模型生态}
\Box{149}{169}{37}{\sffamily\fontsize{8}{11}\selectfont\color{muted}GOSM}
\Box{136}{181}{48}{\sffamily\fontsize{8.5}{12}\selectfont\color{ink}\centering 供给、复用与开放}
\Box{22}{196}{166}{\sffamily\fontsize{8.4}{12}\selectfont\color{muted}\centering 图 1-8　第一章提出方法 · 第二至第五章提供证据 · 第六章归纳洞察与行动}

\Label{22}{207}{166}{贯穿全书的四组研究问题}
\Label{22}{219}{14}{\textcolor{blue}{01}}
\Body{37}{218}{151}{全球协作怎样形成跨层支撑与持续接续的体系能力？}
\Label{22}{231}{14}{\textcolor{teal}{02}}
\Body{37}{230}{151}{中国模型供给如何连接工具链、制度建设与全球社区？}
\Label{22}{243}{14}{\textcolor{gold}{03}}
\Body{37}{242}{151}{全球南方如何通过本地共创和组织支持形成公共能力？}
\Label{22}{255}{14}{\textcolor{violet}{04}}
\Body{37}{254}{151}{中国不同产业、人才群体和城市怎样形成各自优势？}
\Note{22}{270}{166}{第六章依据前文经验证的发现，提炼关键洞察、分层时间前瞻与行动建议；各研究视角互补，不强制一一对应。}
\end{canvas}
\end{document}
